% Einheit 2 – Nr. 30–39
\setcounter{aufgabe}{29}
\einheitenkopf[e2][2 Wachstumsfaktor und Tabelle]{Einheit 2 von 5 · Wachstumsfaktor und Wachstumstabelle}
\zweigzeile{Hier lernst du, aus einem Prozentsatz den Wachstumsfaktor zu bilden und Wachstumstabellen zu ergänzen · neu in diesem Jahr · P10 oft · baut auf: Differenz oder Quotient (Einheit 1), Erhöhen um p\,\%, Quotient zweier Werte}

\begin{aufgabe}{Ich erkenne am Faktor, ob ein Wert zunimmt oder abnimmt.}
Kreuze an. Rechne nichts.
\begin{teile}
\teil Der Wert wird mit 1,4 multipliziert. \quad \kreuz{nimmt zu}\kreuz{nimmt ab}
\teil Der Wert wird mit 0,7 multipliziert. \quad \kreuz{nimmt zu}\kreuz{nimmt ab}
\teil Der Wert wird mit 1,05 multipliziert. \quad \kreuz{nimmt zu}\kreuz{nimmt ab}
\teil Der Wert wird mit 0,95 multipliziert. \quad \kreuz{nimmt zu}\kreuz{nimmt ab}
\teil Der Wert wird mit 1,002 multipliziert. \quad \kreuz{nimmt zu}\kreuz{nimmt ab}
\end{teile}
\end{aufgabe}

\verfahren{Wachstumsfaktor bestimmen}

\begin{aufgabe}{Ich kann aus einem Prozentsatz den Wachstumsfaktor bilden.}
Gib den Wachstumsfaktor an.
\begin{teilezwei}
\tz Zunahme um 10\,\% \quad \feld{q} & \tz Zunahme um 20\,\% \quad \feld{q} \\
\tz Zunahme um 50\,\% \quad \feld{q} & \tz Zunahme um 3\,\% \quad \feld{q} \\
\tz Zunahme um 2,5\,\% \quad \feld{q} & \tz Abnahme um 10\,\% \quad \feld{q} \\
\tz Abnahme um 25\,\% \quad \feld{q} & \tz Abnahme um 4\,\% \quad \feld{q} \\
\tz Abnahme um 0,5\,\% \quad \feld{q} & \tz Zunahme um 100\,\% \quad \feld{q} \\
\end{teilezwei}
\end{aufgabe}

\begin{aufgabe}{Ich kann aus einem Wachstumsfaktor die prozentuale Änderung ablesen.}
Gib an, um wie viel Prozent der Wert zu- oder abnimmt. Schreibe + für Zunahme und $-$ für Abnahme.
\begin{teilezwei}
\tz $q = 1{,}04$ \quad \leerfeld[\%] & \tz $q = 1{,}25$ \quad \leerfeld[\%] \\
\tz $q = 1{,}08$ \quad \leerfeld[\%] & \tz $q = 1{,}6$ \quad \leerfeld[\%] \\
\tz $q = 0{,}9$ \quad \leerfeld[\%] & \tz $q = 0{,}85$ \quad \leerfeld[\%] \\
\tz $q = 1{,}015$ \quad \leerfeld[\%] & \tz $q = 0{,}995$ \quad \leerfeld[\%] \\
\tz $q = 3$ \quad \leerfeld[\%] & \\
\end{teilezwei}
\end{aufgabe}

\begin{aufgabe}{Ich kann den Wachstumsfaktor als Quotient aus einer Tabelle bestimmen.}
Berechne die Quotienten benachbarter Werte und gib den Wachstumsfaktor an.
\begin{teile}
\teil 20; \ 40; \ 80 \quad \feld{q}
\teil 7; \ 21; \ 63 \quad \feld{q}
\teil 300; \ 330; \ 363 \quad \feld{q}
\teil 500; \ 400; \ 320 \quad \feld{q}
\teil 1200; \ 1260; \ 1323; \ 1389,15 \quad Sind alle Quotienten gleich? Antworte in einem Satz und gib $q$ an.
\teil Ein Bestand ist im Jahr 0 genau 400 und im Jahr 2 genau 484. Er wächst jedes Jahr um denselben Prozentsatz. Bestimme den Wachstumsfaktor und den Prozentsatz.
\teil Die Zahl der Bakterien in einer Probe wurde stündlich gezählt:

\sachtabelle{lcccc}{Zeit in h & 0 & 1 & 2 & 3}{Anzahl & 128 & 160 & 200 & 250}

Weise nach, dass die Anzahl jede Stunde mit demselben Faktor wächst, und gib an, um wie viel Prozent sie stündlich zunimmt. (P10 2025 OS)
\end{teile}
\end{aufgabe}

\verfahren{Wachstumstabelle fortschreiben}

\begin{aufgabe}{Ich kann eine Wachstumstabelle mit dem Faktor fortschreiben.}
Ergänze die Tabelle.
\begin{teile}
\teil Wachstumsfaktor 2 \\ \wertetabelle[5,,,]{t}{N}{0,1,2,3}
\teil Wachstumsfaktor 3 \\ \wertetabelle[2,,,]{t}{N}{0,1,2,3}
\teil Wachstumsfaktor 1,1 \\ \wertetabelle[100,,,]{t}{N}{0,1,2,3}
\teil Wachstumsfaktor 1,5 \\ \wertetabelle[40,,,]{t}{N}{0,1,2,3}
\end{teile}
\end{aufgabe}

\begin{aufgabe}{Ich kann eine Wachstumstabelle mit dem Faktor fortschreiben – weiter.}
Ergänze die Tabelle.
\begin{teile}
\teil Ein Bestand von 250 Fischen wächst jedes Jahr um 20\,\%. Trage den Anfangswert und die Werte der nächsten beiden Jahre ein. \\ \wertetabelle[,,]{t}{N}{0,1,2}
\teil Wachstumsfaktor 1,2 \\ \wertetabelle[150,,216,]{t}{N}{0,1,2,3}
\teil Ein Guthaben von 800\,€ wächst jedes Jahr um 5\,\%. Berechne den Wert nach sechs Jahren in einem Schritt. \\ \wertetabelle[800,]{t}{K}{0,6}
\end{teile}
\end{aufgabe}

\begin{aufgabe}{Ich kann eine Wachstumstabelle mit dem Faktor fortschreiben – weiter.}
Ergänze die Tabelle.
\begin{teile}
\teil Wachstumsfaktor 1,1. Berechne den Wert davor. \\ \wertetabelle[,440,484]{t}{N}{0,1,2}
\teil Wachstumsfaktor 1,5. Bestimme die fehlende Zeitangabe und prüfe sie mit dem Faktor. \\ \wertetabelle[128,192,288,972]{t}{N}{0,1,2,{\;}}
\teil Wachstumsfaktor 0,75 \\ \wertetabelle[2400,,,]{t}{N}{0,1,2,3}
\teil Ein Medikament wird im Körper abgebaut, jede Stunde um 15\,\%. Ergänze den fehlenden Wert und die fehlende Zeitangabe. (P10 2016 OS) \\ \wertetabelle[40,34,,{20{,}88}]{t \text{ in h}}{m \text{ in mg}}{0,1,2,{\;}}
\end{teile}
\end{aufgabe}

\begin{aufgabe}{Ich finde den Fehler in einer Wachstumstabelle.}
Finde den Fehler, benenne ihn und korrigiere die Tabelle.
\begin{teile}
\teil Ein Bestand von 500 wächst jährlich um 4\,\%. Lena schreibt: \\ \wertetabelle[500,520,540,560]{t}{N}{0,1,2,3}
\teil Ein Bestand von 200 wächst jährlich um 5\,\%. Kai rechnet mit dem Faktor 0,05: \\ \wertetabelle[200,10,{0{,}5}]{t}{N}{0,1,2}
\teil Ein Guthaben von 800\,€ wächst jährlich um 3\,\%. Jonas schreibt: \\ \wertetabelle[824,{848{,}72}]{t}{K}{0,1}
\end{teile}
\end{aufgabe}

\begin{aufgabe}{Ich kann begründen, woran man den Wachstumsfaktor erkennt.}
Begründe.
\begin{teile}
\teil Warum gehört zu „plus 8\,\%“ der Faktor 1,08 und nicht 0,08?
\teil Ein Bestand hat die Werte 200; \ 220; \ 242; \ 266,2. Tim sagt: „Die Differenzen 20; 22; 24,2 werden größer, also gibt es keinen festen Faktor.“ Hat er recht?
\end{teile}
\end{aufgabe}

\begin{aufgabe}{Ich kann mit einer Wachstumstabelle eine Frage im Sachzusammenhang beantworten.}
Seerosen bedecken 12\,m² eines 50\,m² großen Teichs. Die bedeckte Fläche wächst jede Woche um 25\,\%.
\begin{teile}
\teil Lege eine Tabelle für die Wochen 0 bis 6 an.
\teil In welcher Woche ist zum ersten Mal mehr als die Hälfte des Teichs bedeckt?
\teil Der Besitzer will Seerosen entfernen, bevor 40\,m² bedeckt sind. Bis zu welcher Woche muss er es spätestens tun? Begründe mit der Tabelle.
\end{teile}
\end{aufgabe}
